% Optional related-work insertion near the end of the introduction:
% % Optional related-work insertion near the end of the introduction:
% % Optional related-work insertion near the end of the introduction:
% % Optional related-work insertion near the end of the introduction:
% \input{literature_extras}
% All citation keys below are in references.bib. Keep only the paragraphs
% appropriate to the final paper; the broader generator context is optional.
% Sources checked 2026-09-30 against primary papers/publisher records,
% the formal proof graph, documented intakes and the live Stacks Project.

\paragraph{Related work and proof sources.}
Stafford's original paper proves two-generation of finitely generated torsion
right Weyl-algebra modules and identifies Conjecture~3.8 with their stable
cyclicity~\cite[Theorem~3.7 and pp.~437--438]{Stafford1978}.
Bellamy's 2026 survey records the weaker ordinary-cyclicity question for
finitely generated torsion $A_n(\mathbb C)$-modules as open, and distinguishes
it from Stafford's stronger identity~\cite[Section~3, Conjecture~3.4]{Bellamy2026}.
The machine-checked Stafford38 proof is recorded in Palomar
\cite{Stafford38Palomar2026,Stafford38Zenodo120}.

Constructive predecessors address Stafford's established generator theorems.
Hillebrand--Schmale and Leykin develop effective two-generator methods, and
Leykin also constructs cyclic generators of holonomic modules
\cite{HillebrandSchmale2001,Leykin2004}. Quadrat--Robertz give constructive
module-decomposition and basis algorithms~\cite{QuadratRobertz2007,
QuadratRobertz2014}. Caro--Levcovitz and Caro-Tuesta--Levcovitz study
two-generation and module structure for differential-operator rings with
formal or convergent series coefficients~\cite{CaroLevcovitz2010,
CaroTuestaLevcovitz2020}. These results concern ideal generation, holonomic
cyclicity, and projective-module computation. The constrained identity
$1=dR+FdS$ for arbitrary nonzero $d$ is a further assertion.

The wider generator-theory context includes Stafford's stable-range theorem
and noncommutative Forster--Swan and algebraic $K$-theory results
\cite{StaffordStableRange1981,StaffordGenerating1981}; the latter article
has a corrigendum~\cite{StaffordGeneratingCorrigendum1983} and a related
chapter exposition~\cite{StaffordGeneratingChapter1982}. Stafford also
constructed nontrivial stably free projective right ideals over Weyl
algebras~\cite{StaffordStablyFree1985}. These are surrounding module-structure
results. The OreModules package provides homological and Gr\"obner-basis
methods for multidimensional linear systems over Ore algebras
\cite{OreModules2007} and belongs to the computational background.

For the proof's algebraic framework, Ore's skew-polynomial construction
\cite{Ore1933} supplies historical context for the derivation Ore presentation;
Bj\"ork and Hotta--Takeuchi--Tanisaki develop differential operators, good
filtrations and characteristic varieties~\cite{Bjork,HTT}. The involutivity
input is Gabber's theorem, together with its erratum and the algebraic
formulation of Singh--Kumar~\cite{Gabber1981,Gabber1982,SinghKumar2014}.
Schnell's notes give the noncharacteristic framework
\cite[Lecture~15]{SchnellNotes}. The particular Euler residue, localized
filtered-length contradiction, and asymptotic conormal construction are
developed in this project; the cited background does not attribute those
arguments to the textbooks.

The standard geometric steps can be cited precisely in the Stacks Project:
finite normalization and preservation of projectivity (Lemmas~33.27.1--2,
Tag~0BXQ), the height-one criterion for regularity in a Noetherian normal
domain (Tag~031T), separability of field extensions in characteristic zero
(Tag~0322), and regular functions on proper
geometrically integral varieties (Lemma~33.26.2, Tag~04L2)
\cite[Tags~0BXQ, 031T, 0322, 04L2]{Stacks}.
Normalization is an isomorphism over the smooth locus by the finite
birational normal-target criterion (Tag~0AB1), and completed regular
equicharacteristic local rings have the needed power-series description
(Tag~0C0S)~\cite[Tags~0AB1, 0C0S]{Stacks}. The tangent-lattice limit uses
the valuative criterion~\cite[Chapter~II, Section~4]{hartshorne}.
Faithfully flat descent is covered by Tag~00HP; the recursive finite-order
differential-operator definition and extension to localizations are covered
by Tag~09CH, Lemma~10.133.12~\cite[Tags~00HP, 09CH]{Stacks}.

% Precise edits for the current drafts, to apply during manual polishing:
% human_readable_main.tex: change "every right torsion module" to
% "every finitely generated torsion right A_n(C)-module" when describing
% Bellamy's question. Retain the stronger, arbitrary-field theorem separately.
% Replace \cite{SinghKumar} by \cite{SinghKumar2014}; Bjork is now in the .bib.
% Replace the faithful-flatness \cite{somewhere} by \cite[Tag~00HP]{Stacks}.
% For the other \cite{somewhere}, cite HTT Chapter 1 and give the elementary
% argument: an exhaustive filtration with F_{-1}Q=0 and gr Q=0 has every
% filtration step zero, hence Q=0. A nonzero finite graded module over the
% symbol ring has nonempty support.
% Both paper versions: replace Tag 04L0, Lemma 33.26.2 by Tag 04L2,
% Lemma 33.26.2. Tag 04L0 names the section rather than that lemma.
% Cite Tag 0AB1 at the smooth-locus normalization assertion; Tag 0C0S at
% the completed-local-ring assertion; and Lemma 33.27.2 for projectivity.
% In main.tex's localization section add \cite[Tag~09CH, Lemma~10.133.12]{Stacks}
% at the first definition of finite-order differential operators.
% The exact Euler and filtered-length proofs and the Stafford denominator
% transport remain project arguments, rather than imported textbook results.

% All citation keys below are in references.bib. Keep only the paragraphs
% appropriate to the final paper; the broader generator context is optional.
% Sources checked 2026-09-30 against primary papers/publisher records,
% the formal proof graph, documented intakes and the live Stacks Project.

\paragraph{Related work and proof sources.}
Stafford's original paper proves two-generation of finitely generated torsion
right Weyl-algebra modules and identifies Conjecture~3.8 with their stable
cyclicity~\cite[Theorem~3.7 and pp.~437--438]{Stafford1978}.
Bellamy's 2026 survey records the weaker ordinary-cyclicity question for
finitely generated torsion $A_n(\mathbb C)$-modules as open, and distinguishes
it from Stafford's stronger identity~\cite[Section~3, Conjecture~3.4]{Bellamy2026}.
The machine-checked Stafford38 proof is recorded in Palomar
\cite{Stafford38Palomar2026,Stafford38Zenodo120}.

Constructive predecessors address Stafford's established generator theorems.
Hillebrand--Schmale and Leykin develop effective two-generator methods, and
Leykin also constructs cyclic generators of holonomic modules
\cite{HillebrandSchmale2001,Leykin2004}. Quadrat--Robertz give constructive
module-decomposition and basis algorithms~\cite{QuadratRobertz2007,
QuadratRobertz2014}. Caro--Levcovitz and Caro-Tuesta--Levcovitz study
two-generation and module structure for differential-operator rings with
formal or convergent series coefficients~\cite{CaroLevcovitz2010,
CaroTuestaLevcovitz2020}. These results concern ideal generation, holonomic
cyclicity, and projective-module computation. The constrained identity
$1=dR+FdS$ for arbitrary nonzero $d$ is a further assertion.

The wider generator-theory context includes Stafford's stable-range theorem
and noncommutative Forster--Swan and algebraic $K$-theory results
\cite{StaffordStableRange1981,StaffordGenerating1981}; the latter article
has a corrigendum~\cite{StaffordGeneratingCorrigendum1983} and a related
chapter exposition~\cite{StaffordGeneratingChapter1982}. Stafford also
constructed nontrivial stably free projective right ideals over Weyl
algebras~\cite{StaffordStablyFree1985}. These are surrounding module-structure
results. The OreModules package provides homological and Gr\"obner-basis
methods for multidimensional linear systems over Ore algebras
\cite{OreModules2007} and belongs to the computational background.

For the proof's algebraic framework, Ore's skew-polynomial construction
\cite{Ore1933} supplies historical context for the derivation Ore presentation;
Bj\"ork and Hotta--Takeuchi--Tanisaki develop differential operators, good
filtrations and characteristic varieties~\cite{Bjork,HTT}. The involutivity
input is Gabber's theorem, together with its erratum and the algebraic
formulation of Singh--Kumar~\cite{Gabber1981,Gabber1982,SinghKumar2014}.
Schnell's notes give the noncharacteristic framework
\cite[Lecture~15]{SchnellNotes}. The particular Euler residue, localized
filtered-length contradiction, and asymptotic conormal construction are
developed in this project; the cited background does not attribute those
arguments to the textbooks.

The standard geometric steps can be cited precisely in the Stacks Project:
finite normalization and preservation of projectivity (Lemmas~33.27.1--2,
Tag~0BXQ), the height-one criterion for regularity in a Noetherian normal
domain (Tag~031T), separability of field extensions in characteristic zero
(Tag~0322), and regular functions on proper
geometrically integral varieties (Lemma~33.26.2, Tag~04L2)
\cite[Tags~0BXQ, 031T, 0322, 04L2]{Stacks}.
Normalization is an isomorphism over the smooth locus by the finite
birational normal-target criterion (Tag~0AB1), and completed regular
equicharacteristic local rings have the needed power-series description
(Tag~0C0S)~\cite[Tags~0AB1, 0C0S]{Stacks}. The tangent-lattice limit uses
the valuative criterion~\cite[Chapter~II, Section~4]{hartshorne}.
Faithfully flat descent is covered by Tag~00HP; the recursive finite-order
differential-operator definition and extension to localizations are covered
by Tag~09CH, Lemma~10.133.12~\cite[Tags~00HP, 09CH]{Stacks}.

% Precise edits for the current drafts, to apply during manual polishing:
% human_readable_main.tex: change "every right torsion module" to
% "every finitely generated torsion right A_n(C)-module" when describing
% Bellamy's question. Retain the stronger, arbitrary-field theorem separately.
% Replace \cite{SinghKumar} by \cite{SinghKumar2014}; Bjork is now in the .bib.
% Replace the faithful-flatness \cite{somewhere} by \cite[Tag~00HP]{Stacks}.
% For the other \cite{somewhere}, cite HTT Chapter 1 and give the elementary
% argument: an exhaustive filtration with F_{-1}Q=0 and gr Q=0 has every
% filtration step zero, hence Q=0. A nonzero finite graded module over the
% symbol ring has nonempty support.
% Both paper versions: replace Tag 04L0, Lemma 33.26.2 by Tag 04L2,
% Lemma 33.26.2. Tag 04L0 names the section rather than that lemma.
% Cite Tag 0AB1 at the smooth-locus normalization assertion; Tag 0C0S at
% the completed-local-ring assertion; and Lemma 33.27.2 for projectivity.
% In main.tex's localization section add \cite[Tag~09CH, Lemma~10.133.12]{Stacks}
% at the first definition of finite-order differential operators.
% The exact Euler and filtered-length proofs and the Stafford denominator
% transport remain project arguments, rather than imported textbook results.

% All citation keys below are in references.bib. Keep only the paragraphs
% appropriate to the final paper; the broader generator context is optional.
% Sources checked 2026-09-30 against primary papers/publisher records,
% the formal proof graph, documented intakes and the live Stacks Project.

\paragraph{Related work and proof sources.}
Stafford's original paper proves two-generation of finitely generated torsion
right Weyl-algebra modules and identifies Conjecture~3.8 with their stable
cyclicity~\cite[Theorem~3.7 and pp.~437--438]{Stafford1978}.
Bellamy's 2026 survey records the weaker ordinary-cyclicity question for
finitely generated torsion $A_n(\mathbb C)$-modules as open, and distinguishes
it from Stafford's stronger identity~\cite[Section~3, Conjecture~3.4]{Bellamy2026}.
The machine-checked Stafford38 proof is recorded in Palomar
\cite{Stafford38Palomar2026,Stafford38Zenodo120}.

Constructive predecessors address Stafford's established generator theorems.
Hillebrand--Schmale and Leykin develop effective two-generator methods, and
Leykin also constructs cyclic generators of holonomic modules
\cite{HillebrandSchmale2001,Leykin2004}. Quadrat--Robertz give constructive
module-decomposition and basis algorithms~\cite{QuadratRobertz2007,
QuadratRobertz2014}. Caro--Levcovitz and Caro-Tuesta--Levcovitz study
two-generation and module structure for differential-operator rings with
formal or convergent series coefficients~\cite{CaroLevcovitz2010,
CaroTuestaLevcovitz2020}. These results concern ideal generation, holonomic
cyclicity, and projective-module computation. The constrained identity
$1=dR+FdS$ for arbitrary nonzero $d$ is a further assertion.

The wider generator-theory context includes Stafford's stable-range theorem
and noncommutative Forster--Swan and algebraic $K$-theory results
\cite{StaffordStableRange1981,StaffordGenerating1981}; the latter article
has a corrigendum~\cite{StaffordGeneratingCorrigendum1983} and a related
chapter exposition~\cite{StaffordGeneratingChapter1982}. Stafford also
constructed nontrivial stably free projective right ideals over Weyl
algebras~\cite{StaffordStablyFree1985}. These are surrounding module-structure
results. The OreModules package provides homological and Gr\"obner-basis
methods for multidimensional linear systems over Ore algebras
\cite{OreModules2007} and belongs to the computational background.

For the proof's algebraic framework, Ore's skew-polynomial construction
\cite{Ore1933} supplies historical context for the derivation Ore presentation;
Bj\"ork and Hotta--Takeuchi--Tanisaki develop differential operators, good
filtrations and characteristic varieties~\cite{Bjork,HTT}. The involutivity
input is Gabber's theorem, together with its erratum and the algebraic
formulation of Singh--Kumar~\cite{Gabber1981,Gabber1982,SinghKumar2014}.
Schnell's notes give the noncharacteristic framework
\cite[Lecture~15]{SchnellNotes}. The particular Euler residue, localized
filtered-length contradiction, and asymptotic conormal construction are
developed in this project; the cited background does not attribute those
arguments to the textbooks.

The standard geometric steps can be cited precisely in the Stacks Project:
finite normalization and preservation of projectivity (Lemmas~33.27.1--2,
Tag~0BXQ), the height-one criterion for regularity in a Noetherian normal
domain (Tag~031T), separability of field extensions in characteristic zero
(Tag~0322), and regular functions on proper
geometrically integral varieties (Lemma~33.26.2, Tag~04L2)
\cite[Tags~0BXQ, 031T, 0322, 04L2]{Stacks}.
Normalization is an isomorphism over the smooth locus by the finite
birational normal-target criterion (Tag~0AB1), and completed regular
equicharacteristic local rings have the needed power-series description
(Tag~0C0S)~\cite[Tags~0AB1, 0C0S]{Stacks}. The tangent-lattice limit uses
the valuative criterion~\cite[Chapter~II, Section~4]{hartshorne}.
Faithfully flat descent is covered by Tag~00HP; the recursive finite-order
differential-operator definition and extension to localizations are covered
by Tag~09CH, Lemma~10.133.12~\cite[Tags~00HP, 09CH]{Stacks}.

% Precise edits for the current drafts, to apply during manual polishing:
% human_readable_main.tex: change "every right torsion module" to
% "every finitely generated torsion right A_n(C)-module" when describing
% Bellamy's question. Retain the stronger, arbitrary-field theorem separately.
% Replace \cite{SinghKumar} by \cite{SinghKumar2014}; Bjork is now in the .bib.
% Replace the faithful-flatness \cite{somewhere} by \cite[Tag~00HP]{Stacks}.
% For the other \cite{somewhere}, cite HTT Chapter 1 and give the elementary
% argument: an exhaustive filtration with F_{-1}Q=0 and gr Q=0 has every
% filtration step zero, hence Q=0. A nonzero finite graded module over the
% symbol ring has nonempty support.
% Both paper versions: replace Tag 04L0, Lemma 33.26.2 by Tag 04L2,
% Lemma 33.26.2. Tag 04L0 names the section rather than that lemma.
% Cite Tag 0AB1 at the smooth-locus normalization assertion; Tag 0C0S at
% the completed-local-ring assertion; and Lemma 33.27.2 for projectivity.
% In main.tex's localization section add \cite[Tag~09CH, Lemma~10.133.12]{Stacks}
% at the first definition of finite-order differential operators.
% The exact Euler and filtered-length proofs and the Stafford denominator
% transport remain project arguments, rather than imported textbook results.

% All citation keys below are in references.bib. Keep only the paragraphs
% appropriate to the final paper; the broader generator context is optional.
% Sources checked 2026-09-30 against primary papers/publisher records,
% the formal proof graph, documented intakes and the live Stacks Project.

\paragraph{Related work and proof sources.}
Stafford's original paper proves two-generation of finitely generated torsion
right Weyl-algebra modules and identifies Conjecture~3.8 with their stable
cyclicity~\cite[Theorem~3.7 and pp.~437--438]{Stafford1978}.
Bellamy's 2026 survey records the weaker ordinary-cyclicity question for
finitely generated torsion $A_n(\mathbb C)$-modules as open, and distinguishes
it from Stafford's stronger identity~\cite[Section~3, Conjecture~3.4]{Bellamy2026}.
The machine-checked Stafford38 proof is recorded in Palomar
\cite{Stafford38Palomar2026,Stafford38Zenodo120}.

Constructive predecessors address Stafford's established generator theorems.
Hillebrand--Schmale and Leykin develop effective two-generator methods, and
Leykin also constructs cyclic generators of holonomic modules
\cite{HillebrandSchmale2001,Leykin2004}. Quadrat--Robertz give constructive
module-decomposition and basis algorithms~\cite{QuadratRobertz2007,
QuadratRobertz2014}. Caro--Levcovitz and Caro-Tuesta--Levcovitz study
two-generation and module structure for differential-operator rings with
formal or convergent series coefficients~\cite{CaroLevcovitz2010,
CaroTuestaLevcovitz2020}. These results concern ideal generation, holonomic
cyclicity, and projective-module computation. The constrained identity
$1=dR+FdS$ for arbitrary nonzero $d$ is a further assertion.

The wider generator-theory context includes Stafford's stable-range theorem
and noncommutative Forster--Swan and algebraic $K$-theory results
\cite{StaffordStableRange1981,StaffordGenerating1981}; the latter article
has a corrigendum~\cite{StaffordGeneratingCorrigendum1983} and a related
chapter exposition~\cite{StaffordGeneratingChapter1982}. Stafford also
constructed nontrivial stably free projective right ideals over Weyl
algebras~\cite{StaffordStablyFree1985}. These are surrounding module-structure
results. The OreModules package provides homological and Gr\"obner-basis
methods for multidimensional linear systems over Ore algebras
\cite{OreModules2007} and belongs to the computational background.

For the proof's algebraic framework, Ore's skew-polynomial construction
\cite{Ore1933} supplies historical context for the derivation Ore presentation;
Bj\"ork and Hotta--Takeuchi--Tanisaki develop differential operators, good
filtrations and characteristic varieties~\cite{Bjork,HTT}. The involutivity
input is Gabber's theorem, together with its erratum and the algebraic
formulation of Singh--Kumar~\cite{Gabber1981,Gabber1982,SinghKumar2014}.
Schnell's notes give the noncharacteristic framework
\cite[Lecture~15]{SchnellNotes}. The particular Euler residue, localized
filtered-length contradiction, and asymptotic conormal construction are
developed in this project; the cited background does not attribute those
arguments to the textbooks.

The standard geometric steps can be cited precisely in the Stacks Project:
finite normalization and preservation of projectivity (Lemmas~33.27.1--2,
Tag~0BXQ), the height-one criterion for regularity in a Noetherian normal
domain (Tag~031T), separability of field extensions in characteristic zero
(Tag~0322), and regular functions on proper
geometrically integral varieties (Lemma~33.26.2, Tag~04L2)
\cite[Tags~0BXQ, 031T, 0322, 04L2]{Stacks}.
Normalization is an isomorphism over the smooth locus by the finite
birational normal-target criterion (Tag~0AB1), and completed regular
equicharacteristic local rings have the needed power-series description
(Tag~0C0S)~\cite[Tags~0AB1, 0C0S]{Stacks}. The tangent-lattice limit uses
the valuative criterion~\cite[Chapter~II, Section~4]{hartshorne}.
Faithfully flat descent is covered by Tag~00HP; the recursive finite-order
differential-operator definition and extension to localizations are covered
by Tag~09CH, Lemma~10.133.12~\cite[Tags~00HP, 09CH]{Stacks}.

% Precise edits for the current drafts, to apply during manual polishing:
% human_readable_main.tex: change "every right torsion module" to
% "every finitely generated torsion right A_n(C)-module" when describing
% Bellamy's question. Retain the stronger, arbitrary-field theorem separately.
% Replace \cite{SinghKumar} by \cite{SinghKumar2014}; Bjork is now in the .bib.
% Replace the faithful-flatness \cite{somewhere} by \cite[Tag~00HP]{Stacks}.
% For the other \cite{somewhere}, cite HTT Chapter 1 and give the elementary
% argument: an exhaustive filtration with F_{-1}Q=0 and gr Q=0 has every
% filtration step zero, hence Q=0. A nonzero finite graded module over the
% symbol ring has nonempty support.
% Both paper versions: replace Tag 04L0, Lemma 33.26.2 by Tag 04L2,
% Lemma 33.26.2. Tag 04L0 names the section rather than that lemma.
% Cite Tag 0AB1 at the smooth-locus normalization assertion; Tag 0C0S at
% the completed-local-ring assertion; and Lemma 33.27.2 for projectivity.
% In main.tex's localization section add \cite[Tag~09CH, Lemma~10.133.12]{Stacks}
% at the first definition of finite-order differential operators.
% The exact Euler and filtered-length proofs and the Stafford denominator
% transport remain project arguments, rather than imported textbook results.
